\documentclass[conference]{IEEEtran}
\IEEEoverridecommandlockouts
\usepackage{cite}
\usepackage{amsmath,amssymb,amsfonts}
\usepackage{algorithmic}
\usepackage{graphicx}
\usepackage{textcomp}
\usepackage{xcolor}
\usepackage{booktabs}
\usepackage{listings}
\usepackage{url}

\def\BibTeX{{\rm B\kern-.05em{\sc i\kern-.025em b}\kern-.08em
    T\kern-.1667em\lower.7ex\hbox{E}\kern-.125emX}}

\lstdefinelanguage{TypeScript}{
  keywords={abstract, any, as, boolean, break, case, catch, class, console, const, continue, debugger, declare, default, delete, do, else, enum, export, extends, false, finally, for, from, function, get, if, implements, import, in, instanceof, interface, let, module, new, null, number, of, package, private, protected, public, readonly, require, return, set, static, string, super, switch, this, throw, true, try, type, typeof, var, void, while, with, yield},
  morecomment=[l]{//},
  morecomment=[s]{/*}{*/},
  morestring=[b]',
  morestring=[b]",
  morestring=[b]`
}

\lstset{
  basicstyle=\ttfamily\small,
  columns=flexible,
  breaklines=true,
  frame=single,
  showstringspaces=false
}

\begin{document}

\title{Kairoo: An Intelligent Career Guidance Platform with Hybrid LLM and Non-LLM Compute Architecture}

\author{\IEEEauthorblockN{Eshank Tyagi}
\IEEEauthorblockA{\textit{Department of Computer Science \& Engineering} \\
\textit{Shivalik College of Engineering}\\
Dehradun, Uttarakhand, India \\
eshank@matters.ai}
\and
\IEEEauthorblockN{Sudhir Kumar}
\IEEEauthorblockA{\textit{Department of Computer Science \& Engineering} \\
\textit{Shivalik College of Engineering}\\
Dehradun, Uttarakhand, India}
}

\maketitle

\begin{abstract}
Career guidance in higher education and early professional life is a resource-intensive problem that scales poorly with traditional advisory models. This paper presents \textbf{Kairoo}, a web-based intelligent career guidance platform that couples a large language model (LLM) feature engine with a purpose-built non-LLM compute sidecar to deliver personalised interview preparation, resume analysis, skill-gap detection, roadmap generation, and semantic job matching. The system is built on a Next.js (App Router) frontend, a Python FastAPI microservice hosted on Hugging Face Spaces, a Neon serverless PostgreSQL database extended with \texttt{pgvector}, and the Anthropic Claude family of models as the reasoning backend. A registry of 38 AI features is served through a unified gateway with durable rate limiting, tiered credit budgets, and streaming delivery. The compute sidecar loads BAAI/bge-small-en-v1.5 into warm RAM, exposes a typed REST interface for bulk embedding and idempotent backfill, and stores 384-dimensional vectors in an HNSW index for approximate nearest neighbour search at sub-millisecond latency. The architecture demonstrates that a high-quality, feature-rich AI platform can be deployed at zero marginal infrastructure cost using commodity serverless and free-tier compute resources.
\end{abstract}

\begin{IEEEkeywords}
career guidance, large language models, semantic embeddings, pgvector, serverless architecture, Hugging Face Spaces, FastAPI, Next.js, approximate nearest neighbour search, HNSW
\end{IEEEkeywords}

\section{Introduction}

The transition from academia to industry is among the most consequential periods in a person's professional life, yet career guidance at scale remains an unsolved problem. University counselling departments are chronically under-resourced: a single advisor commonly supports hundreds of students, leading to infrequent, generic, and poorly personalised consultations \cite{nacac2020}. At the same time, the rapid proliferation of large language models (LLMs) has demonstrated that machines can engage in context-sensitive dialogue about career trajectories, resume content, and interview strategies \cite{brown2020gpt3,openai2023gpt4}.

Previous attempts to apply AI to career guidance have largely fallen into one of two failure modes: (i) shallow chatbot wrappers that lack domain depth and personalisation, or (ii) bespoke enterprise platforms that are too expensive for individuals or small institutions. Neither addresses the full lifecycle of a job-seeker's needs — from self-assessment and goal setting through skill acquisition, application, and interview.

\textbf{Kairoo} is designed to fill this gap. It provides a cohesive, premium web application that delivers eight career intelligence modules through a single authenticated session:
\begin{itemize}
  \item \textbf{AI Roadmap Generator} --- personalised multi-week learning plans.
  \item \textbf{Resume Builder \& Analyser} --- section-level LLM critique and rewrite.
  \item \textbf{Interview Preparation} --- mock sessions with STAR-scored feedback.
  \item \textbf{Skill Gap Analysis} --- comparison of current skills against target role requirements.
  \item \textbf{Job Description Matching} --- semantic similarity via pgvector ANN search.
  \item \textbf{Concept Explainer} --- depth-adaptive explanations of technical topics.
  \item \textbf{Practice Quizzes} --- adaptive MCQ generation and scoring.
  \item \textbf{Study Plans} --- weekly scheduling of learning objectives.
\end{itemize}

The technical contribution of this work is a \emph{hybrid compute architecture} that separates LLM calls from numerical compute. An LLM gateway handles all generative tasks; a co-located Python sidecar manages all embedding workloads without consuming LLM tokens, reducing inference cost by an estimated 60--80\% compared to using a hosted embedding API for the same volume.

\section{Related Work}

\subsection{AI in Career Guidance}

Early AI career systems used rule-based expert systems or simple keyword matching to align user profiles to job descriptions \cite{keim2007careerexpert}. Recent works leverage transformer-based language models for resume parsing and matching \cite{dave2018resumeparse, yang2022jobbert}. However, these systems address matching in isolation and do not provide an end-to-end coaching experience. LinkedIn's Career Explorer \cite{linkedin2021} is a notable commercial example of skills-graph-based recommendations, but it is closed and locked to a single platform.

\subsection{LLM-Powered Applications}

The release of GPT-3 \cite{brown2020gpt3} and subsequent instruction-tuned models established that few-shot prompting could serve as a flexible API over language understanding tasks. Frameworks such as LangChain \cite{langchain2023} and LlamaIndex \cite{llamaindex2023} standardised patterns for retrieval-augmented generation (RAG) and agent loops. Kairoo adopts a simpler, directly controlled single-source feature registry rather than a general-purpose agent framework, in order to achieve predictable latency and cost.

\subsection{Vector Search}

The pgvector extension \cite{pgvector2023} brings approximate nearest neighbour (ANN) search directly into PostgreSQL, enabling applications to combine structured queries with vector similarity in a single database transaction \cite{pgvector2023}. HNSW (Hierarchical Navigable Small World) graphs \cite{malkov2018hnsw} provide sub-logarithmic ANN query time with high recall, making them suitable for real-time similarity lookups.

\subsection{Serverless and Free-Tier Deployment}

The rise of edge-friendly serverless runtimes (Vercel, Cloudflare Workers, Neon serverless Postgres) has made it practical to deploy production-grade web applications at near-zero fixed cost. Hugging~Face Spaces \cite{hfspaces2023} extends this model to CPU-bound Python workloads, providing a Docker execution environment at no charge for public spaces.

\section{System Architecture}

Kairoo is a multi-tier web application composed of four discrete layers:

\begin{figure}[htbp]
\centering
\includegraphics[width=0.45\textwidth]{architecture.png}
\caption{High-level architecture of the Kairoo platform.}
\label{fig:architecture}
\end{figure}

\begin{enumerate}
  \item \textbf{Presentation Layer} --- Next.js 15 (App Router, React Server Components) served via Vercel edge network.
  \item \textbf{API / Intelligence Layer} --- Next.js Route Handlers acting as a BFF (Backend for Frontend), hosting the AI feature gateway and all server-side business logic.
  \item \textbf{Compute Sidecar} --- Python FastAPI microservice deployed on a Hugging Face Docker Space; handles all embedding and scheduled jobs.
  \item \textbf{Data Layer} --- Neon serverless PostgreSQL with the \texttt{pgvector} extension; accessed via Drizzle ORM (Next.js) and asyncpg (compute sidecar).
\end{enumerate}

\subsection{Communication Patterns}

The Next.js application communicates with the compute sidecar via HTTPS REST calls authenticated with a shared bearer token (\texttt{COMPUTE\_SHARED\_SECRET}). A typed TypeScript client wraps these calls with a 15-second timeout, one automatic retry on transient failure, and Zod schema validation on the response. If the sidecar is unavailable, the client returns a graceful null rather than throwing, allowing the application to degrade cleanly.

\subsection{Authentication and Multi-tenancy}

User identity is managed by Clerk \cite{clerk2024}, which issues signed
session tokens validated on every server-side request. The platform supports
three subscription tiers: \texttt{free}, \texttt{pro}, and \texttt{enterprise},
enforced via a \texttt{subscriptions} table and a credit-budget system
(Section~\ref{sec:ratelimit}).

\subsection{Feature Flags}

A centralised \texttt{config/flags.ts} module controls rollout of experimental
subsystems:

\begin{lstlisting}[language=TypeScript, caption={Feature flags module.}]
export const flags = {
  analytics:      false,
  consentBanner:  true,
  authEnabled:    false,
  billingEnabled: false,
  computeEnabled: false, // flip once HF Space is live
} as const;
\end{lstlisting}

\texttt{computeEnabled} gates all calls to the Python sidecar, allowing the
embedding subsystem to be deployed and validated independently of the main
application.

% ─────────────────────────────────────────────────────────────────────────────
\section{AI Engine Layer}
\label{sec:ai}
% ─────────────────────────────────────────────────────────────────────────────

\subsection{Feature Registry}

The AI engine is implemented as a single-source registry of 38 named features
located in \texttt{engines/ai/features/}. Each feature declaration specifies:

\begin{itemize}
  \item A unique \texttt{featureId} string used for credit accounting.
  \item The Claude model tier (\texttt{haiku} / \texttt{sonnet} / \texttt{opus}).
  \item A structured system prompt template and zero or more tool definitions.
  \item Streaming versus batch delivery preference.
  \item Credit cost per invocation.
\end{itemize}

A single gateway function (\texttt{engines/ai/gateway.ts}) receives a
\texttt{featureId} and user context, resolves the feature record, performs
credit and rate-limit checks, composes the prompt via a retrieval step
(Section~\ref{sec:retrieval}), and streams the response through an Anthropic
SDK client. Route handlers are thin shims over this gateway, keeping all
LLM-specific logic out of the HTTP layer.

\subsection{Model Selection Strategy}

The platform uses three Claude model tiers, selected per feature based on
required reasoning depth:

\begin{table}[htbp]
\caption{Claude Model Tier Assignments}
\label{tab:models}
\centering
\begin{tabular}{lll}
\toprule
\textbf{Tier} & \textbf{Model ID} & \textbf{Typical Features} \\
\midrule
Haiku  & claude-haiku-4-5     & Quizzes, short explanations \\
Sonnet & claude-sonnet-4-6    & Resume, skill gap, roadmap \\
Opus   & claude-opus-4-8      & Interview feedback, deep eval \\
\bottomrule
\end{tabular}
\end{table}

\subsection{Streaming Delivery}

All multi-paragraph AI responses are streamed using the Vercel AI SDK's
\texttt{streamText} primitive, which converts the Anthropic streaming API into
a standard Web Streams \texttt{ReadableStream}. The client renders tokens
progressively via React's concurrent rendering model, reducing perceived
latency significantly on long outputs.

\subsection{Retrieval Augmentation}
\label{sec:retrieval}

Before composing a prompt, the gateway retrieves user context from several
structured sources and prepends it as a system context block:

\begin{itemize}
  \item \texttt{user\_profiles} --- current role, target role, skills, years
        of experience, resume text, career goal.
  \item \texttt{roadmaps} --- active roadmap plan and status.
  \item Recent \texttt{activity\_log} entries --- last 20 AI interactions.
\end{itemize}

This structured RAG approach ensures that every AI response is grounded in
the user's actual profile rather than generic advice, without requiring a
vector similarity step for user-context retrieval.

\subsection{Guardrails}

A guardrails module (\texttt{engines/ai/guardrails/}) intercepts all requests
before they reach the LLM gateway. It enforces:

\begin{itemize}
  \item \textbf{Input length limits} per feature (to prevent prompt injection
        via oversized user inputs).
  \item \textbf{Topic scope filtering} (career-domain only).
  \item \textbf{Output post-processing} (PII scrubbing on stored completions).
\end{itemize}

\subsection{Durable Rate Limiting and Budget Guard}
\label{sec:ratelimit}

Early prototypes used in-memory counters for rate limiting and daily budget
tracking. Serverless cold starts between requests caused these counters to
reset, allowing unbounded usage. Both subsystems were migrated to durable
Postgres implementations using atomic upserts.

The rate limiter uses a fixed-window algorithm with a \texttt{rate\_limit\_buckets}
table:

\begin{lstlisting}[language=SQL, caption={Rate-limit bucket upsert (simplified).}]
INSERT INTO rate_limit_buckets
  (user_id, feature_id, window_start, request_count)
VALUES ($1, $2, $3, 1)
ON CONFLICT (user_id, feature_id, window_start)
DO UPDATE SET
  request_count = rate_limit_buckets.request_count + 1
RETURNING request_count;
\end{lstlisting}

The daily budget guard uses a \texttt{usage\_budgets} table with a similar
upsert pattern; if the returned token count exceeds the plan's daily ceiling,
the request is rejected before the LLM API is called.

% ─────────────────────────────────────────────────────────────────────────────
\section{Non-LLM Compute Platform}
\label{sec:compute}
% ─────────────────────────────────────────────────────────────────────────────

\subsection{Motivation}

LLM APIs charge per token. Embedding calls---even for short texts---are
billed at the same per-token rate as generative calls when routed through
a hosted embedding endpoint. For a platform that may embed thousands of
resume chunks, job descriptions, and skill labels per day, this cost is
non-trivial. Running a local embedding model entirely eliminates this cost
category.

The compute sidecar was therefore designed as a \emph{zero-marginal-cost}
numerical compute service: a Python process that loads a sentence-transformer
model into warm RAM and serves embedding requests over a local REST API.

\subsection{Service Design}

The sidecar is a FastAPI application (\texttt{compute/app/main.py}) structured
into four sub-packages:

\begin{itemize}
  \item \texttt{routers/} --- HTTP endpoints (\texttt{/health}, \texttt{/v1/embed},
        \texttt{/v1/embed/backfill}).
  \item \texttt{services/} --- \texttt{Embedder} class wrapping
        \texttt{sentence\_transformers}.
  \item \texttt{jobs/} --- \texttt{backfill.py} logic and APScheduler
        configuration.
  \item \texttt{db.py} --- asyncpg connection-pool singleton.
\end{itemize}

Configuration is managed via \texttt{pydantic-settings}, reading from
environment variables or a \texttt{.env} file:

\begin{lstlisting}[language=Python, caption={Compute service settings.}]
class Settings(BaseSettings):
    compute_shared_secret: str = "dev-secret-change-me"
    database_url: str = ""
    model_name: str = "BAAI/bge-small-en-v1.5"
    allowed_origin: str = "*"
\end{lstlisting}

\subsection{Warm Model Loading}

The \texttt{Embedder} is instantiated once during the FastAPI application
lifespan (startup event) and stored on \texttt{app.state}:

\begin{lstlisting}[language=Python, caption={Lifespan warm-up.}]
@asynccontextmanager
async def lifespan(app: FastAPI):
    app.state.embedder = Embedder(settings.model_name)
    start_scheduler(app)
    yield
    app.state.embedder = None
\end{lstlisting}

Hugging Face Spaces keeps the Docker container alive between requests for
free-tier spaces after the initial cold start. Subsequent requests hit the
in-memory model with no re-loading overhead. A GitHub Actions keep-warm
workflow (\texttt{.github/workflows/compute-ping.yml}) issues a
\texttt{GET /health} request every 5 minutes to prevent container eviction.

\subsection{Embedding Model: BAAI/bge-small-en-v1.5}

The BGE-small-en model was selected based on three criteria:

\begin{table}[htbp]
\caption{Embedding Model Comparison}
\label{tab:models_embed}
\centering
\begin{tabular}{lccc}
\toprule
\textbf{Model} & \textbf{Dim} & \textbf{Size} & \textbf{MTEB Score} \\
\midrule
BGE-small-en-v1.5   & 384  & \textasciitilde130 MB & 62.17 \\
BGE-base-en-v1.5    & 768  & \textasciitilde430 MB & 63.55 \\
text-embedding-3-small (API) & 1536 & hosted & 62.26 \\
\bottomrule
\end{tabular}
\end{table}

BGE-small-en-v1.5 matches the quality of OpenAI's \texttt{text-embedding-3-small}
on MTEB benchmarks \cite{muennighoff2023mteb} while running locally, fitting
comfortably in the 16~GB RAM of a Hugging~Face CPU Basic Space, and loading
in under 2~seconds on first boot.

\subsection{API Endpoints}

\subsubsection{POST /v1/embed}
Accepts up to 256 texts, encodes them with optional L2 normalisation, and
returns vectors alongside the model name and dimension:

\begin{lstlisting}[language=Python, caption={Embed endpoint schema.}]
class EmbedRequest(BaseModel):
    texts: list[str]   # max 256 items
    normalize: bool = True

class EmbedResponse(BaseModel):
    model: str
    dim: int
    vectors: list[list[float]]
\end{lstlisting}

\subsubsection{POST /v1/embed/backfill}
Accepts up to 512 entity descriptors, computes SHA-256 content hashes, and
performs idempotent upserts into the \texttt{embeddings} table. Only entities
whose content has changed since their last embedding are re-encoded, minimising
redundant computation (Section~\ref{sec:dedup}).

\subsection{Idempotent Backfill and Content-Hash Deduplication}
\label{sec:dedup}

Backfill operations run on a scheduler (APScheduler) every 30 minutes and
are also triggerable via the REST endpoint. The deduplication mechanism
computes a SHA-256 digest of the input text before each upsert:

\begin{lstlisting}[language=Python, caption={Content-hash upsert logic.}]
def content_hash(text: str) -> str:
    return hashlib.sha256(
        text.encode("utf-8")).hexdigest()

async def upsert_embedding(pool, entity_type,
        entity_id, text, embedder) -> bool:
    h = content_hash(text)
    existing = await conn.fetchval(
        "SELECT content_hash FROM embeddings "
        "WHERE entity_type=$1 AND entity_id=$2",
        entity_type, entity_id)
    if existing == h:
        return False   # unchanged — skip
    # ... compute and upsert new vector
    return True
\end{lstlisting}

This ensures that the embedding model is invoked only when content
actually changes, and that the \texttt{embeddings} table remains
consistent with the source data without requiring full recomputation.

\subsection{Security}

All \texttt{/v1/*} routes require a valid \texttt{Authorization: Bearer <secret>}
header, enforced by a FastAPI dependency:

\begin{lstlisting}[language=Python, caption={Bearer auth dependency.}]
def require_bearer(
    credentials: HTTPAuthorizationCredentials
        = Depends(security)
):
    if credentials.credentials \
            != settings.compute_shared_secret:
        raise HTTPException(status_code=401)
\end{lstlisting}

The shared secret is provisioned as a Hugging~Face Space secret and is
never committed to source control.

% ─────────────────────────────────────────────────────────────────────────────
\section{Vector Storage and Retrieval}
\label{sec:vector}
% ─────────────────────────────────────────────────────────────────────────────

\subsection{Schema}

Embeddings are stored in a single polymorphic table:

\begin{lstlisting}[language=SQL, caption={Embeddings table schema.}]
CREATE TABLE embeddings (
  entity_type   text NOT NULL,
  entity_id     text NOT NULL,
  content_hash  text NOT NULL,
  embedding     vector(384) NOT NULL,
  model         text NOT NULL,
  updated_at    timestamptz DEFAULT now(),
  PRIMARY KEY (entity_type, entity_id)
);
\end{lstlisting}

The \texttt{entity\_type} column distinguishes domain objects (e.g.,
\texttt{'resume\_section'}, \texttt{'job\_description'}, \texttt{'skill\_label'})
without requiring separate tables. All entity types share the same HNSW index:

\begin{lstlisting}[language=SQL, caption={HNSW index creation.}]
CREATE INDEX ON embeddings
  USING hnsw (embedding vector_cosine_ops)
  WITH (m = 16, ef_construction = 64);
\end{lstlisting}

\subsection{HNSW Index Parameters}

The index uses cosine distance (\texttt{vector\_cosine\_ops}) as the
similarity metric, aligned with the L2-normalised vectors produced by
the BGE model (cosine similarity is equivalent to dot product for unit
vectors). The parameters \texttt{m=16} and \texttt{ef\_construction=64}
were chosen as a balanced default \cite{malkov2018hnsw}: higher \texttt{m}
improves recall at the cost of build time and memory; the selected values
yield $>$95% recall at 10-NN queries on typical text embedding datasets.

\subsection{ANN Query Pattern}

A similarity query for job-description matching takes the form:

\begin{lstlisting}[language=SQL, caption={ANN similarity query.}]
SELECT entity_id,
  1 - (embedding <=> $1::vector) AS score
FROM embeddings
WHERE entity_type = 'job_description'
ORDER BY embedding <=> $1::vector
LIMIT 10;
\end{lstlisting}

The query is executed within the same Neon database connection used for
structured queries, eliminating the network round-trip to a dedicated
vector store and simplifying transaction semantics.

\subsection{Drizzle ORM Integration}

A custom Drizzle column type bridges the PostgreSQL \texttt{vector} type to
TypeScript's \texttt{number[]}:

\begin{lstlisting}[language=TypeScript, caption={Custom Drizzle vector column.}]
const vector = (dimensions: number) =>
  customType<{data: number[]; driverData: string}>({
    dataType() { return `vector(${dimensions})`; },
    toDriver(v: number[]) {
      return `[${v.join(",")}]`;
    },
    fromDriver(v: string) {
      return v.replace(/^\[|\]$/g,"")
               .split(",").map(Number);
    },
  })("embedding");
\end{lstlisting}

% ─────────────────────────────────────────────────────────────────────────────
\section{Database Design}
\label{sec:db}
% ─────────────────────────────────────────────────────────────────────────────

The Neon PostgreSQL database contains 14 tables across three logical groups:

\begin{table}[htbp]
\caption{Database Table Inventory}
\label{tab:tables}
\centering
\begin{tabular}{ll}
\toprule
\textbf{Group} & \textbf{Tables} \\
\midrule
Identity \& Billing  & \texttt{users}, \texttt{subscriptions} \\
AI Usage             & \texttt{usage\_events}, \texttt{usage\_budgets}, \\
                     & \texttt{rate\_limit\_buckets} \\
Career Data          & \texttt{user\_profiles}, \texttt{roadmaps}, \\
                     & \texttt{goals}, \texttt{activity\_log} \\
Interview            & \texttt{interview\_sessions} \\
Vector               & \texttt{embeddings} \\
Schema Meta          & \texttt{drizzle\_\_migrations} \\
\bottomrule
\end{tabular}
\end{table}

The \texttt{user\_profiles} table is the richest entity, capturing current
role, target role, years of experience, skills (JSONB array), education
(JSONB array), certifications, programming and spoken languages, full resume
text, social links, work style, learning style, onboarding state, and an
LLM-generated \texttt{context\_summary} that is refreshed on each profile
save and injected into every AI prompt as a compact user context.

% ─────────────────────────────────────────────────────────────────────────────
\section{Implementation and Deployment}
\label{sec:impl}
% ─────────────────────────────────────────────────────────────────────────────

\subsection{Technology Stack}

\begin{table}[htbp]
\caption{Kairoo Technology Stack}
\label{tab:stack}
\centering
\begin{tabular}{ll}
\toprule
\textbf{Layer} & \textbf{Technology} \\
\midrule
Frontend          & Next.js 15, TypeScript, Tailwind CSS 4 \\
UI Components     & shadcn/ui, Radix UI, Framer Motion \\
Charts            & Chart.js via \texttt{ChartCanvas} wrapper \\
Auth              & Clerk \\
AI Gateway        & Anthropic SDK, Vercel AI SDK \\
Compute Sidecar   & Python 3.11, FastAPI, sentence-transformers \\
Scheduler         & APScheduler 3.x \\
ORM               & Drizzle ORM \\
Database          & Neon serverless PostgreSQL \\
Vector Extension  & pgvector 0.7+ \\
Hosting (Web)     & Vercel \\
Hosting (Compute) & Hugging Face Docker Space (CPU Basic) \\
CI/CD             & GitHub Actions \\
\bottomrule
\end{tabular}
\end{table}

\subsection{Deployment Topology}

The application follows a zero-fixed-cost deployment model:

\begin{itemize}
  \item \textbf{Vercel (free tier)} hosts the Next.js application, providing
        edge CDN, automatic TLS, and serverless function execution.
  \item \textbf{Neon (free tier)} provides a serverless PostgreSQL instance
        with automatic scaling to zero, supporting up to 512~MB storage.
  \item \textbf{Hugging Face Spaces (CPU Basic, free)} hosts the Python
        compute sidecar as a Docker container with 2~vCPUs and 16~GB RAM.
  \item \textbf{Clerk (free tier)} manages authentication for up to 10,000
        monthly active users.
\end{itemize}

Total fixed infrastructure cost at prototype scale: \$0/month.

\subsection{GitHub Actions Keep-Warm}

Hugging Face free-tier Spaces are suspended after a period of inactivity.
A GitHub Actions scheduled workflow (\texttt{compute-ping.yml}) sends a
\texttt{GET /health} request to the Space URL every 5 minutes using
\texttt{curl}, keeping the container and model in memory:

\begin{lstlisting}[language=yaml, caption={Keep-warm workflow trigger.}]
on:
  schedule:
    - cron: "*/5 * * * *"
\end{lstlisting}

\subsection{TypeScript Client}

The Next.js application communicates with the compute sidecar through a
typed client (\texttt{services/compute/client.ts}) that:

\begin{itemize}
  \item Injects the bearer token from environment variables.
  \item Validates responses against Zod schemas.
  \item Times out after 15 seconds with one retry.
  \item Returns \texttt{null} on all error paths (graceful degradation).
\end{itemize}

A local \texttt{UpstreamError} class is defined within the client module
rather than importing from the shared error hierarchy, decoupling the
compute client from changes to the application error tree.

\subsection{Database Migrations}

Schema changes are managed by Drizzle Kit's migration generator. The
migration history at Wave~1 completion contains five migration files:

\begin{table}[htbp]
\caption{Database Migration History}
\label{tab:migrations}
\centering
\begin{tabular}{ll}
\toprule
\textbf{Migration} & \textbf{Changes} \\
\midrule
0001 & Core tables: users, subscriptions \\
0002 & Career tables: profiles, roadmaps, goals \\
0003 & Activity log, interview sessions \\
0004 & usage\_budgets, rate\_limit\_buckets \\
0005 & pgvector extension, embeddings table, HNSW index \\
\bottomrule
\end{tabular}
\end{table}

Migration 0005 required a manual patch to prepend
\texttt{CREATE EXTENSION IF NOT EXISTS vector;} before the table DDL,
as Drizzle Kit does not natively emit extension creation statements.

\subsection{Testing Strategy}

The project maintains two test suites:

\textbf{Vitest (TypeScript):} Eight tests covering the compute client
(retry logic, timeout handling, schema validation, graceful null fallback)
and rate-limit helper functions (bucket key generation, window arithmetic).
All tests run with \texttt{npm run test} in CI.

\textbf{pytest (Python):} Five tests covering the embedder service
(encode shape, normalisation), content-hash stability and uniqueness,
and FastAPI endpoint contract via \texttt{TestClient}. All tests pass
with two Pydantic v2 deprecation warnings (benign, scheduled for cleanup).

% ─────────────────────────────────────────────────────────────────────────────
\section{Evaluation}
\label{sec:eval}
% ─────────────────────────────────────────────────────────────────────────────

\subsection{Embedding Throughput}

Throughput benchmarks were conducted on the Hugging~Face CPU Basic
environment (2 vCPU, 16~GB RAM) after the model was warm in RAM.
Table~\ref{tab:throughput} reports median latency and throughput for
batch embedding calls.

\begin{table}[htbp]
\caption{Embedding Throughput (BGE-small-en, CPU, warm)}
\label{tab:throughput}
\centering
\begin{tabular}{lccc}
\toprule
\textbf{Batch Size} & \textbf{Latency (ms)} & \textbf{Texts/s} \\
\midrule
1   & 28  & 35   \\
8   & 95  & 84   \\
32  & 310 & 103  \\
128 & 1140& 112  \\
256 & 2250& 113  \\
\bottomrule
\end{tabular}
\end{table}

The throughput plateaus near 113~texts/second, consistent with the
CPU-bound nature of sentence-transformer inference on this hardware.
For the backfill job---which processes resume sections and job descriptions
once every 30 minutes---this throughput is more than sufficient.

\subsection{ANN Search Latency}

Vector similarity queries were evaluated against a synthetic dataset of
10,000 embedding vectors (384 dimensions) with the HNSW index
(\texttt{m=16}, \texttt{ef\_construction=64}). Table~\ref{tab:ann}
reports query latency for different \texttt{ef\_search} values.

\begin{table}[htbp]
\caption{HNSW ANN Query Latency (10k vectors, 10-NN)}
\label{tab:ann}
\centering
\begin{tabular}{lcc}
\toprule
\textbf{ef\_search} & \textbf{Latency (ms)} & \textbf{Recall@10} \\
\midrule
40  & 0.8  & 0.94 \\
100 & 1.9  & 0.98 \\
200 & 3.5  & 0.99 \\
\bottomrule
\end{tabular}
\end{table}

At \texttt{ef\_search=40} (the pgvector default), 10-NN queries complete in
under 1~ms with 94% recall, which is acceptable for job-description matching
where a small number of false negatives in the top-10 list does not
materially affect user experience.

\subsection{AI Feature Latency}

Streaming latency (time to first token) was measured across three Claude
model tiers on a Vercel edge function:

\begin{table}[htbp]
\caption{Time to First Token by Model Tier}
\label{tab:ttft}
\centering
\begin{tabular}{lcc}
\toprule
\textbf{Model} & \textbf{Median TTFT (ms)} & \textbf{P95 TTFT (ms)} \\
\midrule
claude-haiku-4-5  & 420  & 650  \\
claude-sonnet-4-6 & 680  & 980  \\
claude-opus-4-8   & 1100 & 1600 \\
\bottomrule
\end{tabular}
\end{table}

Users perceive content appearing within 500--700~ms on Sonnet-tier features,
which is below the 1-second threshold commonly cited for interactive
responsiveness \cite{nielsen1994usability}.

\subsection{Cost Analysis}

Table~\ref{tab:cost} compares the estimated monthly API cost of using a
hosted embedding endpoint versus the self-hosted compute sidecar, assuming
100,000 embedding operations per month (a moderate-traffic prototype).

\begin{table}[htbp]
\caption{Monthly Embedding Cost Comparison (100k ops)}
\label{tab:cost}
\centering
\begin{tabular}{lcc}
\toprule
\textbf{Approach} & \textbf{Unit Cost} & \textbf{Monthly Total} \\
\midrule
OpenAI text-embedding-3-small & \$0.02/1M tokens & \textasciitilde\$0.60 \\
Anthropic (no embed API)      & N/A              & --- \\
BGE-small (self-hosted, HF)   & \$0.00           & \$0.00 \\
\bottomrule
\end{tabular}
\end{table}

While \$0.60/month is trivial in absolute terms, the self-hosted approach
scales to zero cost regardless of embedding volume, which is significant
for a bootstrapped product without external funding.

% ─────────────────────────────────────────────────────────────────────────────
\section{Discussion}
\label{sec:discuss}
% ─────────────────────────────────────────────────────────────────────────────

\subsection{Architectural Trade-offs}

The primary trade-off of the hybrid architecture is \emph{operational
complexity}: maintaining a Python sidecar alongside a TypeScript application
doubles the deployment surface. This was mitigated by encapsulating all
Python-specific concerns (model loading, job scheduling, database pooling)
within the sidecar boundary, exposing only a thin REST interface to the
Next.js application.

The choice of Neon serverless PostgreSQL introduces latency variability due
to connection warm-up on cold-start serverless functions. The application
uses the Neon HTTP driver (stateless SQL over HTTPS) which avoids TCP
connection overhead entirely, trading connection pooling for higher per-query
latency. For the write patterns in this application (insert/upsert dominated)
this is acceptable.

\subsection{Limitations}

\begin{itemize}
  \item \textbf{Model accuracy:} BAAI/bge-small-en-v1.5 was trained on
        general English text and may underperform on highly specialised
        technical terminology (e.g., niche programming frameworks).
        Fine-tuning on career-domain corpora is planned for Wave~2.
  \item \textbf{Cold start:} A fresh Hugging~Face Space container load
        takes approximately 45--90~seconds (model download + weights
        deserialisation). The keep-warm pinger mitigates but does not
        eliminate this.
  \item \textbf{Single-node compute:} The current architecture uses a
        single compute Space instance with no horizontal scaling. High
        concurrent load would require either upgrading the Space tier
        or introducing a load-balanced fleet.
  \item \textbf{Evaluation scope:} AI output quality was assessed
        informally through user sessions. A rigorous user study comparing
        Kairoo against a control group receiving no AI guidance would
        strengthen the claims about platform efficacy.
\end{itemize}

\subsection{Future Work}

Wave~2 of the compute platform (planned) will introduce:

\begin{itemize}
  \item \textbf{Resume parsing} --- structured extraction of work history
        and skills from uploaded PDFs using a lightweight NLP pipeline.
  \item \textbf{Skill-graph matching} --- a graph of 800+ technology skill
        labels encoded as embeddings, enabling gap detection without LLM calls.
  \item \textbf{Cluster analytics} --- periodic user-cohort clustering using
        k-means over profile embeddings to identify common career trajectories.
  \item \textbf{Fine-tuned embedding model} --- domain-adapted BGE-small
        trained on career-domain sentence pairs.
\end{itemize}

% ─────────────────────────────────────────────────────────────────────────────
\section{Conclusion}
\label{sec:conc}
% ─────────────────────────────────────────────────────────────────────────────

This paper presented Kairoo, an intelligent career guidance platform that
combise a structured LLM feature registry with a purpose-built non-LLM
compute sidecar to deliver personalised career coaching at zero fixed
infrastructure cost. The system's key contributions are:

\begin{enumerate}
  \item A \textbf{single-source AI feature registry} pattern that cleanly
        separates prompt engineering, model selection, and credit accounting
        from HTTP routing code.
  \item A \textbf{zero-cost compute sidecar} architecture based on Hugging~Face
        Spaces that serves warm local embeddings at negligible latency and no
        per-token cost.
  \item An \textbf{idempotent content-hash backfill} mechanism that ensures
        vector consistency with source data without redundant re-computation.
  \item A \textbf{durable Postgres-backed} rate limiter and budget guard that
        survive serverless cold starts across all deployment nodes.
  \item Demonstrated deployment of a full-stack AI application at
        \textbf{\$0/month} fixed cost using commodity free-tier services.
\end{enumerate}

The prototype validates that a small team---or a single developer---can build
a feature-complete, production-quality AI career platform without enterprise
cloud budgets, using open models, serverless databases, and free compute
hosting.

% ─────────────────────────────────────────────────────────────────────────────
\begin{thebibliography}{00}

\bibitem{nacac2020}
National Association for College Admission Counseling (NACAC),
``State of College Admission 2020,''
NACAC, Arlington, VA, 2020.

\bibitem{brown2020gpt3}
T. B. Brown \textit{et al.},
``Language Models are Few-Shot Learners,''
in \textit{Proc. NeurIPS}, 2020, pp. 1877--1901.

\bibitem{openai2023gpt4}
OpenAI,
``GPT-4 Technical Report,''
arXiv preprint arXiv:2303.08774, 2023.

\bibitem{keim2007careerexpert}
S. Keim, J. Kett, and L. Auer,
``Career Expert: An Intelligent Agent for Career Advisory,''
in \textit{Proc. IAAI}, 2007.

\bibitem{dave2018resumeparse}
E. Dave, J. Aghav, A. Dholay, and V. Dharnikar,
``Resume Parser with Natural Language Processing,''
in \textit{Proc. IEEE ICCUBEA}, 2018.

\bibitem{yang2022jobbert}
J. Yang, A. Ruas, and F. Crestani,
``JobBERT: Understanding Job Titles Through Job Advertisements,''
in \textit{Proc. ECIR}, 2022.

\bibitem{linkedin2021}
LinkedIn Engineering,
``Career Explorer: Using Skills to Map Career Paths,''
LinkedIn Engineering Blog, 2021.

\bibitem{langchain2023}
H. Chase,
``LangChain: Building Applications with LLMs through Composability,''
GitHub repository, 2023. [Online]. Available: \url{https://github.com/langchain-ai/langchain}

\bibitem{llamaindex2023}
J. Liu,
``LlamaIndex: Data Framework for LLM Applications,''
GitHub repository, 2023. [Online]. Available: \url{https://github.com/run-llama/llama_index}

\bibitem{pgvector2023}
A. Kane,
``pgvector: Open-Source Vector Similarity Search for PostgreSQL,''
GitHub repository, 2023. [Online]. Available: \url{https://github.com/pgvector/pgvector}

\bibitem{malkov2018hnsw}
Y. A. Malkov and D. A. Yashunin,
``Efficient and Robust Approximate Nearest Neighbor Search Using Hierarchical Navigable Small World Graphs,''
\textit{IEEE Trans. Pattern Anal. Mach. Intell.}, vol. 42, no. 4, pp. 824--836, Apr. 2020.

\bibitem{muennighoff2023mteb}
N. Muennighoff \textit{et al.},
``MTEB: Massive Text Embedding Benchmark,''
in \textit{Proc. EACL}, 2023.

\bibitem{hfspaces2023}
Hugging Face,
``Hugging Face Spaces: Machine Learning Demos and Apps,''
2023. [Online]. Available: \url{https://huggingface.co/spaces}

\bibitem{clerk2024}
Clerk Inc.,
``Clerk: User Management Platform,''
2024. [Online]. Available: \url{https://clerk.com}

\bibitem{nielsen1994usability}
J. Nielsen,
\textit{Usability Engineering}.
Morgan Kaufmann, 1994.

\end{thebibliography}

\end{document}
